\scanpage{26}

\begin{center}
ALGEBRE COMMUTATIVE

CHAPITRE IV - DECOMPOSITION PRIMAIRE

\uline{Décomposition primaire et extension des scalaires} (Etat 2)

-:-:-
\end{center}
\margin{Rédaction n°355}

L'état 1 porte le n°352. Après l'avoir lu (et même avant), Grothendieck
a signalé au rédacteur qu'on pouvait généraliser le résultat principal, en
remplaçant l'un des anneaux par un module. La démonstration s'étend sans
changement, et c'est bien plus joli. Le rédacteur en a profité pour enlever
quelques hypothèses noethériennes inutiles.

\begin{center}-:-:-\end{center}

Dans tout ce qui suit, on désigne par $A$ et $B$ deux anneaux commutatifs,
et l'on considère un homomorphisme $\pi:A\longrightarrow B$. Cet homomorphisme fait de
$B$ une $A$-algèbre. On se donne en outre un $B$-module $F$ que l'on suppose \uline{plat}
\uline{sur} $A$.

PROPOSITION 1.- \uline{Soit} $\underline p$ \uline{un idéal premier de} $A$. \uline{Pour tout} $\underline P\in\operatorname{Ass}_B(F/\underline pF)$ \uline{on a}
$\pi^{-1}(\underline P)=\underline p$.

Puisque $\pi(\underline p)$ annule $F/\underline pF$, on a $\pi(\underline p)\subset\underline P$, d'où $\underline p\subset\pi^{-1}(\underline P)$. D'autre
part, si $a\in A-\underline p$, l'homothétie de rapport $a$ est injective dans $A/\underline p$, donc
aussi (puisque $F$ est $A$-plat) dans $F/\underline pF$, ce qui montre que $\pi(a)\notin\underline P$, CQFD.

PROPOSITION 2.- \uline{Soit} $E$ \uline{un} $A$-\uline{module}. \uline{On a} :
\[
(\ast)\qquad
\operatorname{Ass}_B(E\otimes_A F)\supset\bigcup_{\underline p\in\operatorname{Ass}_A(E)}\operatorname{Ass}_B(F/\underline pF)
\]
\uline{Lorsque} $A$ \uline{est noethérien}, \uline{l'inclusion} $(\ast)$ \uline{est une égalité}.

Soit $\underline p\in\operatorname{Ass}_A(E)$. Par définition de $\operatorname{Ass}_A(E)$, il existe une suite exacte :
\diagram{\begin{tikzcd}
0\arrow[r]&A/\underline p\arrow[r]&E.
\end{tikzcd}}

Puisque $F$ est plat sur $A$, on en déduit une suite exacte :
\diagram{\begin{tikzcd}
0\arrow[r]&F/\underline pF\arrow[r]&E\otimes_A F,
\end{tikzcd}}

\margin{25}
